%% ============================================================
%%  Appendix A: Additional Experimental Results
%% ============================================================
\chapter{Additional Experimental Results}
\label{app:additional_results}

% Use the appendix for figures and tables that support the main
% text but would break the flow of Chapter 4 if included inline.
% Good candidates: per-seed breakdowns, full results tables,
% supplementary sensitivity plots.


\section{Full Results Table}
\label{app:full_table}

Full per-run result data (all methods, maps, seeds, communication ranges,
and task loads) are available in the accompanying code repository.
A condensed cross-map results summary is provided in Chapter~\ref{sec:exp_throughput}
for all six experimental maps, integrating both stable-regime and stress-regime
steady-state results.


\section{Rerun Interval Sensitivity}
\label{app:ri_sensitivity}

Dedicated rerun-interval sensitivity experiments were run for
$ri \in \{10, 25, 50, 100, 200\}$ timesteps on all six maps.
Preliminary analysis of the results indicates that completion rate
and throughput are largely insensitive to $ri$ above 25 at
light-to-moderate loads; $ri = 10$ increases allocation overhead without
meaningful throughput gains; $ri = 200$ introduces latency in responding
to topology changes at high loads.
The canonical $ri = 50$ used throughout the thesis sits near the
flat region of this sensitivity curve.


\section{Stress-Regime Timeout Diagnostics}
\label{app:stress_timeout_diagnostics}

The stress-regime timeout diagnostics are moved out of the main chapter
because they primarily confirm the failure boundary already described in
Chapter~\ref{sec:exp_throughput}. The underlying image files remain in their
current locations on disk; only the LaTeX inclusion location changes.

\begin{figure}[htbp]
	\centering
	\begin{subfigure}[t]{0.32\textwidth}
		\centering
		\includegraphics[width=\textwidth]{figures/gridworld_kiva/20260426_195056/ss_results/compact_views/compact_heatmap_timeout_rate.png}
		\caption{Kiva}
	\end{subfigure}
	\hfill
	\begin{subfigure}[t]{0.32\textwidth}
		\centering
		\includegraphics[width=\textwidth]{figures/gridworld_kiva_large/20260427_050329/ss_results/compact_views/compact_heatmap_timeout_rate.png}
		\caption{Kiva-Large}
	\end{subfigure}
	\hfill
	\begin{subfigure}[t]{0.32\textwidth}
		\centering
		\includegraphics[width=\textwidth]{figures/gridworld_shelf_aisle/20260427_105146/ss_results/compact_views/compact_heatmap_timeout_rate.png}
		\caption{Shelf-Aisle}
	\end{subfigure}
	\caption{Stress-regime timeout-rate heatmaps moved from the main chapter. The panels document the failure boundary but add limited information beyond the accompanying throughput and capacity-frontier plots.}
	\label{fig:appendix_stress_timeout_heatmaps}
\end{figure}


\section{Per-Seed Breakdown}
\label{app:per_seed}

For all main-analysis maps and load points, results are computed from two
random seeds (42 and 123).
Seed variance is low in the stable operating regime: for \texttt{warehouse\_small}
at $\comm \geq 15$ and $\text{tc} \leq 96$, makespans across seeds differ
by fewer than 20\% in completed runs, and all key conclusions (LCBA
outperforms CBBA at low comm range, sub-optimality ratio within $2\times$
of SGA) hold on both seeds independently.
Variance is higher at $\comm = 4$--$8$ and high load, consistent with
the greater sensitivity to initial allocation in disconnected regimes.


\section{Graph Connectivity Analysis}
\label{app:graph_connectivity}

The relationship between communication graph structure and task completion
is visible in the completion-degradation curves in
Figure~\ref{fig:batch_degradation_main_maps}: completion rate remains stable
until a threshold communication range is crossed, below which it drops
sharply for CBBA but only gradually for LCBA.
